\chapter{The Particle Spectrum from Codec Geometry}

In the previous chapter we proved four exact results from the single
definition $B = \rho - \mathrm{diag}(\rho)$: Pauli exclusion, the
universal boson, the Born rule identity, and the virtual propagator
lifecycle. All followed from the codec structure alone, with no
interaction terms introduced.

This chapter extends the analysis. The previous results concerned a
fermion hopping between two sites. The question now is: what happens
when a fermion is delocalised across $n$ sites? The answer will be
three exact theorems about the norm structure of any fermion, a
conservation law that subsumes all previous results, and a
classification of the particle spectrum that requires no gauge
symmetry groups, coupling constants, or postulated particle species.

\section{The Universal Norm Theorems}

Consider a fermion in equal superposition across $n$ sites:

\[
\psi_k = \frac{1}{\sqrt{n}}\,e^{i\phi_k}, \quad k = 0, \ldots, n-1,
\]

for arbitrary real phases $\phi_k$. Split the density matrix
$\rho = |\psi\rangle\langle\psi|$ into its diagonal and off-diagonal
parts:

\[
F = \mathrm{diag}(\rho), \qquad B = \rho - \mathrm{diag}(\rho).
\]

Three things are true about these objects, for any $n$ and any choice
of phases.

\textbf{Universal Fermion Norm.} $\|F\| = 1/\sqrt{n}$.

\textit{Proof.} $F_k = |\psi_k|^2 = 1/n$ for all $k$, so $\|F\|^2 =
\sum_k 1/n^2 = 1/n$. \qed

\textbf{Universal Boson Norm.} $\|B\| = \sqrt{(n-1)/n}$.

\textit{Proof.} $|\rho_{ij}|^2 = 1/n^2$ for all $i \neq j$. There
are $n(n-1)$ off-diagonal pairs, giving $\|B\|^2 = n(n-1)/n^2 =
(n-1)/n$. \qed

\textbf{Conservation Law.} $\|F\|^2 + \|B\|^2 = 1$.

\textit{Proof.} For any pure state $\rho^2 = \rho$, so
$\|\rho\|^2 = \mathrm{Tr}(\rho^2) = \mathrm{Tr}(\rho) = 1$. Since
$F$ and $B$ occupy orthogonal subspaces, the result follows. \qed

All three results are phase-independent. They hold for any $n \geq 1$
and any choice of $\phi_k$. The conservation law subsumes the
two-site result of the previous chapter: setting $n = 2$ and
$\psi(\theta) = \cos\theta\,e_i + i\sin\theta\,e_j$ recovers
$\|F\|^2 + \|B\|^2 = 1$ as a special case.

The physical reading is direct. As a fermion delocalises across more
sites, its fermionic content $\|F\|$ decreases as $1/\sqrt{n}$ while
its bosonic content $\|B\|$ increases toward 1. The more spread out
the fermion, the richer its compression residual. The conservation
law ensures these always balance.

\begin{table}[h]
\centering
\begin{tabular}{ccccc}
\toprule
$n$ & $\|F\|$ & $\|B\|$ & $\|F\|^2$ & $\|B\|^2$ \\
\midrule
1 & $1$          & $0$          & $1$   & $0$ \\
2 & $1/\sqrt{2}$ & $1/\sqrt{2}$ & $1/2$ & $1/2$ \\
3 & $1/\sqrt{3}$ & $\sqrt{2/3}$ & $1/3$ & $2/3$ \\
4 & $1/2$        & $\sqrt{3}/2$ & $1/4$ & $3/4$ \\
\bottomrule
\end{tabular}
\caption{Universal fermion and boson norms. Every row satisfies
$\|F\|^2 + \|B\|^2 = 1$ exactly.}
\end{table}

\section{The Winding Number and Internal Symmetry}

The norms $\|F\|$ and $\|B\|$ depend only on $n$. The internal
structure of the compression residual --- its symmetry under
transposition --- depends on the phase pattern, specifically on the
winding number $m$: the number of full phase rotations the pattern
makes around the ring of $n$ sites.

The natural minimum-complexity phase patterns on an $n$-site ring are
the discrete Fourier modes $\phi_k^{(m)} = 2\pi mk/n$, ordered by
ascending spectral complexity. The symmetry of $B$ under these
patterns determines whether the fermion can exist as a free asymptotic
state:

\begin{itemize}
    \item \textbf{Antisymmetric} ($B = -B^\top$, achieved by $m$ odd
    for $n = 2$): the residual has self-contained antisymmetric
    structure. The fermion can propagate as a free asymptotic state.

    \item \textbf{Symmetric} ($B = B^\top$, achieved by $m = 0$):
    the residual has self-contained symmetric structure and propagates
    freely with scalar character.

    \item \textbf{Mixed symmetry} ($n \geq 3$, $m \geq 1$): the
    residual has no self-contained symmetry class. Projecting onto
    either the symmetric or antisymmetric subspace leaves a nonzero
    residual. The fermion cannot be described as a free asymptotic
    state without reference to its source configuration.
\end{itemize}

The pair $(n, m)$ --- sites and winding number --- defines a
\emph{fermion complexity class}. This is the complete specification.
No gauge groups, coupling constants, or symmetry postulates are
required.

\section{The $n = 2$ Sector: Leptons}

The two-site fermion has $\|F\| = \|B\| = 1/\sqrt{2}$ --- equal
fermionic and bosonic content.

For winding $m = 1$, the compression residual $B$ is purely
antisymmetric, with eigenvalues $\pm 1/2$ and the universal $-\pi/2$
phase of the previous chapter. The fermion has a self-contained
antisymmetric residual and propagates freely. This is the charged
lepton --- electron, muon, or tau, depending on the complexity level.

What Standard Model language calls a photon is the name given to the
correlation between two such fermion events. When one charged lepton
recoils, another does so in a correlated way. The codec encodes this
correlation in the off-diagonal of $\rho$. No separate photon particle
is required.

For winding $m = 0$, the residual $B$ is purely symmetric, with zero
net phase winding. The fermion has no electromagnetic phase structure
and is electrically neutral. This is the neutrino.

Electric charge, within the $n = 2$ sector, is the winding number.
$m = 0$ gives a neutral fermion; $m = 1$ gives a charged fermion. No
additional postulate is needed.

\section{The $n = 3$ Sector: Quarks and Confinement}

For three-site equal superposition, each site carries probability
$|\psi_k|^2 = 1/3$. This is the per-site occupancy of a quark in a
colour triplet. The value $1/3$ is not postulated; it is the Born-rule
probability of finding the fermion on any one colour site.

For winding $m = 1$, the compression residual has off-diagonal entries
of equal magnitude:

\[
|B_{ij}| = \frac{1}{3} \quad \text{for all colour pairs } i \neq j.
\]

This is exact colour symmetry, emerging from the equal-superposition
constraint and the universal boson norm theorem. No $\mathrm{SU}(3)$
symmetry is imposed.

What Standard Model language calls a gluon is the correlation between
quark colour changes. The mixed symmetry of $B$ for $n = 3$, $m = 1$
means this correlation has no self-contained symmetry class. Projecting
onto either self-contained class leaves a nonzero residual of magnitude
approximately $0.408$, verified numerically. A mixed-symmetry fermion
always requires reference to the configuration from which it arose.

This is the information-theoretic basis of confinement. A fermion whose
compression residual has mixed symmetry cannot be described as a free
asymptotic state under any minimum-description-length codec. Colour
confinement is a structural consequence of three-site codec geometry,
not a separately imposed force law. This remains a conjecture --- the
open problem is to prove it as a theorem.

\section{The Complete Classification}

The fermion complexity classes organise the particle spectrum without
remainder. The table below lists the exact results and distinguishes
them from conjectures.

\begin{table}[h]
\centering
\begin{tabular}{cclll}
\toprule
$(n,m)$ & Symmetry & Fermion class & SM label & Status \\
\midrule
$(1,0)$ & --- & Localised, no residual & vacuum & exact \\
$(2,0)$ & symmetric & Neutral lepton & neutrino & exact \\
$(2,1)$ & antisymmetric & Charged lepton & electron/muon/tau & exact \\
$(3,0)$ & symmetric & Colour-neutral scalar & Higgs sector & conjecture \\
$(3,1)$ & mixed & Colour-charged, confined & quark/gluon & exact \\
$(4,2)$ & symmetric & Spin-2 class & open & open \\
\bottomrule
\end{tabular}
\caption{Fermion complexity classes. The SM label column gives the
Standard Model name for the correlation pattern an observer would
attribute to this class --- derived labels, not fundamental particles.}
\end{table}

The SM labels in the right column are not names of fundamental
objects. They are the names given to patterns of fermion behaviour by
observers who do not know they are watching a compression codec. No
detector has ever registered a photon, gluon, or $W$ boson as a
localised pixel. Detectors register fermion events. The force carriers
are the story told about the correlations between those events.

\section{The Mass Hierarchy}

Under Solomonoff induction, $P(\psi) \propto 2^{-C_s(\psi)}$, so
observable fermions are ordered by ascending spectral complexity.
The minimum fermionic codec unit is 4 bits, arising from the four
binary choices required to specify one fermion hop: which site, which
frame, real or imaginary component, and sign of the phase.

The three charged lepton generations are the same $n = 2$, $m = 1$
fermion class at successive complexity levels:

\begin{table}[h]
\centering
\begin{tabular}{lrrl}
\toprule
Fermion & Mass (MeV) & $\log_2(m/m_e)$ & Nearest integer \\
\midrule
Electron & $0.511$    & $0.00$  & $0$ \\
Muon     & $105.66$   & $7.69$  & $8$ \\
Tau      & $1776.86$  & $11.76$ & $12$ \\
\bottomrule
\end{tabular}
\caption{Charged lepton masses in $\log_2(m/m_e)$ space. The
near-integer sequence with 4-bit granularity suggests a discrete
complexity ladder.}
\end{table}

The separations are approximately $0$, $8$, and $12$ bits relative to
the electron --- near-integer multiples of the 4-bit codec unit. The
deviations from exact integers are interpreted as compression gains
under Solomonoff induction.

A falsifiable prediction follows: no stable charged fermion should
exist beyond the tau, because the lognormal probability distribution
of Chapter~3 places a finite upper bound on observable complexity. No
fourth charged lepton has been observed.

The quark mass hierarchy spans several orders of magnitude and has
not yet been placed on the complexity ladder. This is the primary open
calculation.

\section{Two Layers and the Bridge Between Them}

The framework now has two independent derivations running in parallel.

Chapters~2 through 5 derived the geometric structure of spacetime ---
expansion, gravity, black hole entropy, the Friedmann equation ---
from the aspect ratio and information conservation, without reference
to the wavefunction codec.

Chapters~6 through 8 derive quantum mechanics and the particle
spectrum from the wavefunction codec, without reference to the
geometric layer.

Both derivations are self-consistent and exact within their domain.
The question is how they connect.

The natural bridge is the worldline. A fermion's trajectory through
spacetime is a complex worldline $\psi_{\mathrm{traj}}(t) = x(t) +
iy(t)$. The minimum spectral complexity closed worldline is a single
Fourier mode --- a complex exponential tracing an ellipse. This is
Kepler's orbit. The minimum-$C_s$ worldline in the geometric spacetime
is exactly the trajectory selected by the fermion's Solomonoff
induction.

The geometric and codec layers are not independent. They are two
projections of the same compression principle onto different physical
degrees of freedom. The derivation of the bridge --- showing that the
Einstein field equations emerge as the stationarity condition of $C_s$
over metric configurations --- is the primary target of the next
chapter.

\section{What This Chapter Establishes}

Three exact theorems and one conservation law follow from the
definition $B = \rho - \mathrm{diag}(\rho)$ applied to $n$-site
equal superpositions: universal fermion norm $1/\sqrt{n}$, universal
boson norm $\sqrt{(n-1)/n}$, and conservation $\|F\|^2 + \|B\|^2 =
1$ for any pure state.

The pair $(n, m)$ classifies all fermion complexity classes. The
$n = 2$ sector gives charged leptons and neutrinos distinguished by
winding number. The $n = 3$ sector gives confined colour-charged
fermions and a scalar sector. The mass hierarchy organises into a
near-integer sequence with 4-bit granularity and predicts no stable
charged fermion beyond the tau.

Three open problems remain: why three lepton generations and not four,
why the quark masses have the values they have, and whether the $1/r^2$
electromagnetic force law follows from the $n = 2$, $m = 1$ fermion
residual by the same flux argument that gives gravity. These are
calculation problems, not conceptual ones. The framework is already
in place.
